\documentclass[11pt,a4paper]{article}

% ── Packages ──────────────────────────────────────────────────────────────────
\usepackage[top=2.8cm, bottom=2.8cm, left=3cm, right=3cm]{geometry}
\usepackage[T1]{fontenc}
\usepackage[utf8]{inputenc}
\usepackage{lmodern}
\usepackage{microtype}
\usepackage[colorlinks=true, linkcolor=linkblue, urlcolor=accent,
            pdftitle={L1-05 Databases Relational -- System Design Foundations},
            bookmarks=true]{hyperref}
\usepackage{xcolor}
\usepackage{titlesec}
\usepackage{enumitem}
\usepackage{booktabs}
\usepackage{tabularx}
\usepackage{tcolorbox}
\usepackage{fancyhdr}
\usepackage{amsmath}
\usepackage{parskip}
\usepackage{mdframed}
\usepackage{graphicx}
\usepackage{tikz}
\usepackage{setspace}
\usepackage{soul}
\usepackage{array}
\usepackage{longtable}

% ── Colors ────────────────────────────────────────────────────────────────────
\definecolor{primary}{HTML}{1A1A2E}
\definecolor{accent}{HTML}{C0392B}
\definecolor{highlight}{HTML}{1A5276}
\definecolor{linkblue}{HTML}{154360}
\definecolor{lightbg}{HTML}{F4F6F7}
\definecolor{quizbg}{HTML}{EBF5FB}
\definecolor{termbg}{HTML}{EAF2FF}
\definecolor{curiobg}{HTML}{F5EEF8}
\definecolor{greenbg}{HTML}{EAFAF1}
\definecolor{notebg}{HTML}{FDF2E9}
\definecolor{accentlight}{HTML}{FADBD8}

% ── Section styling ───────────────────────────────────────────────────────────
\titleformat{\section}
  {\large\bfseries\color{highlight}}
  {\thesection.}{0.6em}{}
  [\vspace{2pt}\color{highlight!60}\rule{\linewidth}{0.6pt}]

\titleformat{\subsection}
  {\normalsize\bfseries\color{primary}}
  {\thesubsection}{0.5em}{}

\titleformat{\subsubsection}
  {\normalsize\itshape\color{primary!80}}
  {}{0em}{}

\titlespacing{\section}{0pt}{18pt}{8pt}
\titlespacing{\subsection}{0pt}{12pt}{4pt}
\titlespacing{\subsubsection}{0pt}{8pt}{2pt}

% ── Header/Footer ─────────────────────────────────────────────────────────────
\pagestyle{fancy}
\fancyhf{}
\fancyhead[L]{\small\color{highlight}\textbf{L1-05 --- Databases: Relational}}
\fancyhead[R]{\small\color{primary!70}System Design Interview Prep}
\fancyfoot[C]{\small\color{primary!60}\thepage}
\renewcommand{\headrulewidth}{0.4pt}
\renewcommand{\headrule}{\hbox to\headwidth{\color{highlight!40}\leaders\hrule height \headrulewidth\hfill}}

% ── Box environments ──────────────────────────────────────────────────────────
\tcbuselibrary{skins, breakable, hooks}

\newtcolorbox{termbox}[1]{
  enhanced, breakable,
  colback=termbg, colframe=highlight, coltitle=white,
  fonttitle=\bfseries\small, title={#1},
  left=8pt, right=8pt, top=6pt, bottom=6pt,
  attach boxed title to top left={yshift=-2.5mm, xshift=6mm},
  boxed title style={colback=highlight, sharp corners, left=4pt, right=4pt}
}

\newtcolorbox{industrybox}[1]{
  enhanced, breakable,
  colback=greenbg, colframe=highlight!50!black, coltitle=white,
  fonttitle=\bfseries\small, title={Industry Data: #1},
  left=8pt, right=8pt, top=6pt, bottom=6pt,
  attach boxed title to top left={yshift=-2.5mm, xshift=6mm},
  boxed title style={colback=highlight!60!black, sharp corners, left=4pt, right=4pt}
}

\newtcolorbox{trapbox}[1]{
  enhanced,
  colback=accentlight, colframe=accent, coltitle=white,
  fonttitle=\bfseries\small, title={Interview Trap: #1},
  left=8pt, right=8pt, top=6pt, bottom=6pt,
  attach boxed title to top left={yshift=-2.5mm, xshift=6mm},
  boxed title style={colback=accent, sharp corners, left=4pt, right=4pt}
}

\newtcolorbox{thinkbox}{
  enhanced,
  colback=notebg, colframe=accent!60,
  left=8pt, right=8pt, top=6pt, bottom=6pt,
  leftrule=3pt, rightrule=0pt, toprule=0pt, bottomrule=0pt
}

\newtcolorbox{curiobox}[1]{
  enhanced, breakable,
  colback=curiobg, colframe=primary!60, coltitle=white,
  fonttitle=\bfseries\small, title={Q: #1},
  left=8pt, right=8pt, top=6pt, bottom=6pt,
  attach boxed title to top left={yshift=-2.5mm, xshift=6mm},
  boxed title style={colback=primary!70, sharp corners, left=4pt, right=4pt}
}

\newtcolorbox{quizbox}{
  enhanced, breakable,
  colback=quizbg, colframe=highlight, coltitle=white,
  fonttitle=\bfseries, title={Self-Quiz},
  left=10pt, right=10pt, top=8pt, bottom=8pt,
  attach boxed title to top left={yshift=-2.5mm, xshift=6mm},
  boxed title style={colback=highlight, sharp corners, left=4pt, right=4pt}
}

\newtcolorbox{answerbox}{
  enhanced, breakable,
  colback=lightbg, colframe=primary!40, coltitle=white,
  fonttitle=\bfseries\small, title={Model Answers},
  left=10pt, right=10pt, top=8pt, bottom=8pt,
  attach boxed title to top left={yshift=-2.5mm, xshift=6mm},
  boxed title style={colback=primary!60, sharp corners, left=4pt, right=4pt}
}

\newtcolorbox{insightbox}{
  enhanced,
  colback=primary!5, colframe=primary,
  left=8pt, right=8pt, top=6pt, bottom=6pt,
  leftrule=4pt, rightrule=0.4pt, toprule=0.4pt, bottomrule=0.4pt
}

\setstretch{1.15}

% ─────────────────────────────────────────────────────────────────────────────
\begin{document}

% ── Title Page ───────────────────────────────────────────────────────────────
\thispagestyle{empty}
\begin{center}
  \vspace*{2.5cm}
  {\color{highlight}\rule{\linewidth}{1.5pt}}
  \vspace{0.6cm}

  {\small\color{primary!60}\textbf{SYSTEM DESIGN INTERVIEW PREP} \quad|\quad Layer 1: Foundations}

  \vspace{0.4cm}
  {\fontsize{32}{40}\selectfont\bfseries\color{primary} L1-05}

  \vspace{0.3cm}
  {\fontsize{22}{28}\selectfont\bfseries\color{highlight} Databases: Relational}

  \vspace{0.2cm}
  {\large\color{primary!70}\itshape ACID, Indexes, Replication, and Sharding}

  \vspace{0.5cm}
  {\color{highlight}\rule{0.4\linewidth}{0.6pt}}

  \vspace{0.6cm}
  {\small\color{primary!60}
    Edition: 2025 \quad|\quad Data sources: PostgreSQL documentation, AWS Aurora blog,\\
    Shopify Engineering Blog, Notion Engineering, PlanetScale architecture docs,\\
    Instagram Engineering archive, Percona blog
  }

  \vspace{1.2cm}
  \begin{insightbox}
    \textbf{Why this lesson is assumed knowledge in every interview.}

    Every system design interview that involves persistent data -- which is nearly all of them -- will eventually ask: ``What database would you use?'' If you say ``PostgreSQL'' without being able to explain what that actually buys you (ACID transactions, B-tree indexes, row-level locking via MVCC), you have named a brand, not made an engineering decision. Interviewers at Apple ICT3 and equivalent levels do not need to hear that you know PostgreSQL exists; they need to hear you reason about \emph{why} its consistency guarantees make it the right choice for a payments table, and why those same guarantees become a bottleneck at 50 million writes per day. This lesson gives you the internal model to have that conversation fluently, without hesitation.
  \end{insightbox}

  \vspace{1cm}
  {\color{highlight}\rule{\linewidth}{1.5pt}}
\end{center}

\newpage

% ── Table of Contents ─────────────────────────────────────────────────────────
\tableofcontents

\newpage

% =============================================================================
\section{Why This Exists: The Problem Before Relational Databases}
% =============================================================================

Imagine you are building a payroll system in 1970. You store employee records in flat text files on disk -- one file per department, one line per employee. This works until the day your application crashes halfway through processing a raise: forty records are updated, sixty are not, and you have no reliable way to know which is which. You re-run the job, but now some employees have received two raises and others none. Restoring the state requires reading backup tapes, replaying manual logs, and a lot of embarrassed phone calls to the finance department.

This is not a contrived scenario. It describes exactly the class of problem that Edgar Codd was observing in enterprise computing when he published ``A Relational Model of Data for Large Shared Data Banks'' in 1970. The core insight was simple but revolutionary: data should be organized into relations (tables) with mathematically defined rules governing how records relate to one another, and the database engine itself -- not the application -- should be responsible for enforcing those rules. The application developer should never need to write code that worries about partial updates, inconsistent reads by concurrent users, or what happens when the power fails mid-write.

The relational model gave us three things that flat files could not: a declarative query language (SQL) that separated \emph{what} you want from \emph{how} to find it, referential integrity constraints that made it impossible to create an orphaned order for a non-existent customer, and transaction semantics that made the partial-update problem disappear. These three properties -- not any specific implementation -- are what ``relational database'' actually means.

IBM built System R, the first prototype relational database, in the mid-1970s. Oracle shipped a commercial version in 1979. PostgreSQL traces its lineage directly to the POSTGRES project at UC Berkeley, started by Michael Stonebraker in 1986. MySQL emerged as an open-source alternative in 1995. By the time the web era exploded in the 2000s, the relational database was so ubiquitous that it was simply called ``the database.'' Every junior engineer learned SQL before they learned anything else. And this is precisely why interviewers assume you know it: it is the default, and deviating from it requires justification.

\newpage

% =============================================================================
\section{ACID Properties: The Contract a Relational Database Makes}
% =============================================================================

The word ACID is thrown around in interviews almost as a password -- candidates say it to signal they know what they are talking about, without always being able to explain what it means in concrete terms. This section gives you the concrete terms.

\subsection{Atomicity: All or Nothing}

\begin{termbox}{Atomicity}
A transaction is treated as a single indivisible unit. Either \emph{all} operations within the transaction succeed and are committed, or \emph{none} of them are. There is no visible intermediate state.
\end{termbox}

The payroll example from the previous section is precisely an atomicity failure. Atomicity is implemented in modern databases through a write-ahead log (WAL). Before any data page on disk is modified, the database first writes a description of the intended change to a log file. If the system crashes mid-transaction, the database replays the WAL on recovery and either completes or rolls back the transaction to a clean state. The data pages on disk never contain half-committed state because the log is always written first.

\begin{thinkbox}
\textbf{But what if the WAL write itself crashes?} The WAL is written sequentially, which is fast, and the \texttt{fsync()} system call forces the log to durable storage before the transaction is acknowledged as committed. If the system crashes before \texttt{fsync()} completes, the transaction is simply not in the log and is treated as if it never happened. The database does not acknowledge the commit to the client until after the \texttt{fsync()} returns successfully.
\end{thinkbox}

\subsection{Consistency: Integrity Constraints Are Always Satisfied}

\begin{termbox}{Consistency}
A transaction brings the database from one valid state to another. ``Valid'' means all declared constraints -- primary keys, foreign keys, unique constraints, check constraints -- remain satisfied before and after the transaction.
\end{termbox}

Consistency is the only ACID property that is partly the application's responsibility. The database enforces the constraints you declare; it cannot enforce business rules you never told it about. If your schema allows a negative account balance because you forgot to add a \texttt{CHECK (balance >= 0)} constraint, the database will happily store negative balances. The ``C'' in ACID is therefore a shared contract: the database guarantees mechanical integrity; the developer guarantees correct schema design.

\subsection{Isolation: Concurrent Transactions Do Not Interfere}

Isolation is the most nuanced ACID property and the one most likely to surface in a senior engineering interview. The question it answers is: what does one transaction see when another transaction is running concurrently?

\begin{termbox}{Isolation Levels (weakest to strongest)}
SQL defines four standard isolation levels, each trading performance for consistency:

\begin{enumerate}[leftmargin=*, itemsep=2pt]
  \item \textbf{Read Uncommitted} -- a transaction can read data written by another transaction that has not yet committed. Almost never used in practice because it allows ``dirty reads.''
  \item \textbf{Read Committed} -- a transaction only reads data that has been committed. This is the default in PostgreSQL and Oracle. Prevents dirty reads but allows ``non-repeatable reads'' (the same row read twice in a transaction can return different values if another transaction commits between the two reads).
  \item \textbf{Repeatable Read} -- within a transaction, re-reading the same row always returns the same value, even if another transaction commits a change in between. Prevents non-repeatable reads. MySQL's InnoDB uses this as its default. Can still allow ``phantom reads'' (new rows inserted by another transaction may appear in range queries).
  \item \textbf{Serializable} -- the strongest level. Transactions appear to execute in some serial order, as if they ran one after another with no overlap. Prevents dirty reads, non-repeatable reads, and phantom reads. The most expensive because it typically requires range locks or predicate locking.
\end{enumerate}
\end{termbox}

\begin{industrybox}{PostgreSQL -- MVCC Avoids Read Locks}
PostgreSQL implements isolation through Multi-Version Concurrency Control (MVCC). Rather than locking rows when a writer modifies them, PostgreSQL keeps multiple versions of each row simultaneously. A reader always sees a consistent snapshot of the database as of the moment its transaction began -- without taking any locks. Writers and readers therefore never block each other. This is why PostgreSQL can sustain high concurrent read-write workloads without readers queuing behind writers. The tradeoff is ``dead tuples'': old row versions accumulate and must be periodically cleaned up by a background process called \texttt{VACUUM}. At Notion's scale (2024), their engineering team documented that poorly tuned autovacuum was one of the primary sources of PostgreSQL performance degradation as their user base grew into tens of millions.
\end{industrybox}

\subsection{Durability: Committed Data Survives Crashes}

\begin{termbox}{Durability}
Once a transaction is committed, the changes persist even if the system crashes immediately afterward. The database guarantees this by ensuring committed data has been written to durable storage (disk or equivalent) before acknowledging the commit.
\end{termbox}

Durability is implemented through the same WAL mechanism that enables atomicity. When a transaction commits, PostgreSQL calls \texttt{fsync()} on the WAL segment. This forces the operating system's page cache to flush to physical disk. On modern SSDs with capacitors (enterprise-grade), this is extremely reliable. On cheap virtual machines with no hardware write cache, \texttt{fsync()} can be the biggest bottleneck in write throughput.

\begin{trapbox}{``NoSQL is faster because it sacrifices ACID''}
This is one of the most common misconceptions in system design interviews. Many NoSQL databases -- including MongoDB (since 4.0, 2018) and CockroachDB -- support ACID transactions. The performance difference between relational and NoSQL systems is not about ACID; it is about data model flexibility, horizontal scaling architecture, and query patterns. Saying ``I'll use MongoDB because ACID is too slow'' signals a misunderstanding. The correct framing is: ``For this access pattern, a document model avoids multi-table joins and reduces round trips, which improves latency at the application layer.''
\end{trapbox}

\newpage

% =============================================================================
\section{B-Tree Indexes: Why Queries Are Fast}
% =============================================================================

A table with ten million rows and no index requires the database to scan every row to find the ones matching your \texttt{WHERE} clause. This is called a sequential scan. For a table on disk where each page holds 100 rows, finding a single record means reading up to 100,000 pages. At typical SSD read speeds, this is tens of milliseconds for a single lookup -- unacceptable for a web application.

Indexes are a separate data structure maintained alongside the table that makes lookups fast. The most common index structure in relational databases -- used by PostgreSQL, MySQL InnoDB, Oracle, and SQL Server -- is the B-tree.

\subsection{How B-Trees Work}

A B-tree (balanced tree) stores index keys in sorted order across a tree of fixed-size nodes called pages. Each internal node holds $k$ to $2k$ keys and $k+1$ to $2k+1$ child pointers, keeping the tree balanced. The leaf nodes of the tree contain the actual index values along with pointers (row IDs, also called tuple IDs or TIDs) to the actual data rows on the heap.

When you query \texttt{WHERE user\_id = 42}, the database traverses the B-tree from the root. At each level it performs a binary search within the node to find the child pointer to follow. After traversing $O(\log n)$ levels (where $n$ is the number of rows), it reaches a leaf node and reads the row pointer, then fetches the actual data row from the heap. For a table of ten million rows, a B-tree of height 4 or 5 suffices, reducing 100,000 page reads to 5 or 6.

\begin{thinkbox}
\textbf{Why not a hash table, which is $O(1)$?} PostgreSQL also supports hash indexes, but they have a critical limitation: they only support equality lookups (\texttt{=}). A B-tree supports range queries (\texttt{BETWEEN}, \texttt{<}, \texttt{>}), \texttt{ORDER BY}, and \texttt{LIKE 'prefix\%'} because the keys are stored in sorted order. Range queries are extremely common in real applications (``find all orders placed in the last 7 days''), which is why B-trees dominate in practice.
\end{thinkbox}

\subsection{Composite Indexes and the Leftmost Prefix Rule}

A composite index covers multiple columns. An index on \texttt{(user\_id, created\_at)} stores rows sorted first by \texttt{user\_id}, then by \texttt{created\_at} within each user. This index can answer queries that filter on \texttt{user\_id} alone (it can use the left portion of the sort key), queries that filter on \texttt{(user\_id, created\_at)} together (full use), and range scans on \texttt{created\_at} \emph{within a specific user\_id}. It cannot efficiently answer queries that filter on \texttt{created\_at} alone, because the rows are not sorted globally by that column -- they are sorted by \texttt{user\_id} first.

This is the \textbf{leftmost prefix rule}: an index on \texttt{(A, B, C)} can be used for queries that filter on \texttt{A}, on \texttt{(A, B)}, or on \texttt{(A, B, C)} -- but not for queries that filter only on \texttt{B} or only on \texttt{C}.

\subsection{Covering Indexes: Eliminating Heap Fetches}

A covering index contains all the columns a query needs, so the database can answer the query entirely from the index without touching the heap (the main table data). For example, if your query is \texttt{SELECT user\_id, email FROM users WHERE user\_id = 42}, an index on \texttt{(user\_id, email)} is a covering index for this query -- no heap access needed.

\begin{industrybox}{Shopify -- Index Strategy at Scale}
Shopify runs one of the largest MySQL deployments in the world, processing millions of orders per day across tens of thousands of merchant storefronts. Their engineering blog (2023) documents a recurring pattern: as tables grow past 500 million rows, queries that were fast with covering indexes begin degrading because the index itself becomes too large to fit in the buffer pool (MySQL's in-memory page cache). At that scale, Shopify found that decomposing composite indexes and restructuring queries to use smaller, more selective leading columns was more effective than simply adding memory. Index size discipline -- keeping indexes narrow and selective -- becomes as important as having indexes at all.
\end{industrybox}

\subsection{Index Selectivity and When Not to Index}

An index on a column with very few distinct values (``gender'' with two values, or a boolean flag) has poor selectivity. If half the rows match your filter, scanning the index and then fetching 50\% of the heap is slower than a sequential scan -- the optimizer will likely ignore the index and do a full table scan instead. High selectivity (many distinct values relative to total rows) makes an index valuable. Low selectivity makes it overhead without benefit.

Every index also carries a write cost: every \texttt{INSERT}, \texttt{UPDATE}, or \texttt{DELETE} must update all indexes on that table. A table with fifteen indexes pays fifteen times the write overhead compared to a table with none. Over-indexing is a real performance problem, not a theoretical one.

\newpage

% =============================================================================
\section{Transactions, Locking, and Concurrency Control}
% =============================================================================

Transactions are the mechanism by which ACID guarantees are applied to groups of statements. Understanding how databases implement concurrency control -- and what goes wrong when you get it wrong -- is the difference between a candidate who knows SQL and one who can design production systems.

\subsection{Row-Level Locking in PostgreSQL}

PostgreSQL implements locking at two levels. Most \texttt{SELECT} statements take no locks at all (thanks to MVCC), while \texttt{INSERT}, \texttt{UPDATE}, and \texttt{DELETE} statements take row-level locks on the specific rows they modify. Table-level locks are only taken for operations that restructure the table itself: \texttt{ALTER TABLE}, \texttt{TRUNCATE}, or explicit \texttt{LOCK TABLE} commands.

Row-level locking means two transactions can concurrently update different rows in the same table without blocking each other. This is fundamentally different from table-level locking (which PostgreSQL uses for older versions or explicit overrides), where any write to a table blocks all other writes to that table. MySQL's older MyISAM storage engine used table-level locking, which is why InnoDB replaced it as the default storage engine in MySQL 5.5 (2010).

\begin{termbox}{Deadlock}
A deadlock occurs when transaction A holds a lock on row 1 and is waiting for a lock on row 2, while transaction B holds a lock on row 2 and is waiting for a lock on row 1. Neither can proceed. PostgreSQL detects this cycle and automatically rolls back one of the transactions, returning an error to the application. The application must retry the rolled-back transaction. A well-designed application always handles deadlock errors, because deadlocks are inevitable in concurrent systems.
\end{termbox}

\subsection{Optimistic vs. Pessimistic Concurrency}

The MVCC approach in PostgreSQL is a form of \textbf{optimistic concurrency control}: transactions proceed without taking read locks, assuming conflicts are rare. At commit time, the database checks whether any data the transaction read has been modified by another committed transaction; if so, it aborts and the application retries.

\textbf{Pessimistic concurrency control} explicitly locks rows before reading them, using \texttt{SELECT ... FOR UPDATE}. This prevents other transactions from modifying the locked rows until the lock holder commits or rolls back. Pessimistic locking is the right choice when you know conflicts are likely -- for example, decrementing a shared inventory counter -- because it avoids the retry overhead of repeated optimistic conflicts. It is the wrong choice for high-read, low-conflict workloads, because it turns reads into blocking operations.

\newpage

% =============================================================================
\section{Replication: Scaling Beyond a Single Machine}
% =============================================================================

A single database server has a ceiling: one CPU, one set of disks, one network card. Every system that grows past a certain size must therefore answer the question of how to spread the data and the load across multiple machines. Replication is the first and most common answer.

\subsection{Primary-Replica Replication}

\begin{termbox}{Primary-Replica (Leader-Follower) Replication}
One server is designated the \emph{primary} (also called master or leader). All writes go to the primary. The primary streams its WAL to one or more \emph{replica} servers (also called secondaries, followers, or read replicas), which apply the same operations in order. At any point in time, a replica is a lagging but consistent copy of the primary.
\end{termbox}

The write path is simple: the application sends writes to the primary, the primary writes to its WAL, commits, and streams WAL records to replicas. The primary does not wait for replicas to acknowledge by default -- this is \textbf{asynchronous replication}. The replicas apply WAL records as fast as they arrive, but there is always some lag.

\begin{termbox}{Replication Lag}
The delay between when a write is committed on the primary and when it becomes visible on a replica. On a healthy, lightly loaded replica, this is typically under 100 milliseconds. Under heavy write load, or when a replica is catching up after maintenance, lag can grow to seconds or minutes.
\end{termbox}

\begin{thinkbox}
\textbf{What happens if you route a read to a replica immediately after writing?} You might read stale data -- the write has not replicated yet. This is called \emph{read-after-write inconsistency}. The solution is to route reads-of-your-own-writes to the primary (using sticky sessions or a read-your-writes token), and let all other reads go to replicas. This is exactly what Instagram's engineering team described in their database architecture posts: critical reads (e.g., ``did my post just publish?'') hit the primary; feed queries hit replicas.
\end{thinkbox}

\subsection{Synchronous vs. Asynchronous Replication}

In synchronous replication, the primary waits for at least one replica to acknowledge that it has durably written the WAL record before acknowledging the write to the application. This eliminates replication lag for the acknowledged replica but increases write latency by one network round trip. PostgreSQL supports this via the \texttt{synchronous\_standby\_names} setting.

\begin{industrybox}{AWS Aurora -- Six Copies Across Three Availability Zones}
Amazon Aurora, AWS's cloud-native relational database, takes a radically different approach to replication. Rather than replicating the entire WAL across TCP connections, Aurora separates storage from compute. The storage layer maintains six copies of every 10GB segment of the database, spread across three availability zones in pairs (two copies per AZ). Writes are committed once 4 of 6 storage nodes acknowledge -- a quorum write. The compute layer (the Aurora instance running PostgreSQL or MySQL-compatible SQL) never writes data files directly; it only writes log records to the storage layer. Aurora's 2023 engineering blog documents write latency of under 1 ms for the storage acknowledgment on average, because the storage nodes are colocated in the same AWS region and connected via dedicated network fabric. A read replica in Aurora shares the same storage volume as the primary -- it does not need to receive or apply WAL; it simply reads from the shared storage. This allows Aurora read replicas to have near-zero replication lag compared to traditional PostgreSQL streaming replication.
\end{industrybox}

\subsection{Failover and High Availability}

When the primary goes down, a replica must be promoted to become the new primary. In asynchronous replication, any writes that were committed on the primary but not yet replicated are permanently lost (this is the replication lag window turned into data loss). In synchronous replication, the promoted replica has everything the primary had (up to the last acknowledged write), so there is no data loss, but the system was slower during normal operation.

Tools like Patroni (for PostgreSQL) and Orchestrator (for MySQL) automate failover: they detect primary failure, elect a new primary from the replicas, and update DNS or a virtual IP so that the application automatically connects to the new primary. This is how Shopify, GitHub, and most large PostgreSQL users handle HA (high availability) without manual intervention.

\newpage

% =============================================================================
\section{Sharding: Horizontal Partitioning for Write Scale}
% =============================================================================

Read replicas solve the \emph{read} scaling problem. But what about writes? All writes still go to a single primary. A busy e-commerce database during Black Friday -- say, Shopify processing 40,000 requests per second at peak (documented 2023) -- cannot funnel all writes through one PostgreSQL primary indefinitely. At some point, the single primary becomes the bottleneck, and no amount of hardware upgrade changes that. This is where sharding enters.

\subsection{What Sharding Is}

\begin{termbox}{Sharding (Horizontal Partitioning)}
Splitting a large dataset across multiple independent database instances (shards), each responsible for a non-overlapping subset of the data. Each shard is a full database with its own primary and replicas. Queries that touch a single shard are routed directly. Queries that touch multiple shards must be fanned out and the results aggregated.
\end{termbox}

Sharding is not a feature built into vanilla PostgreSQL or MySQL. It is an architectural pattern implemented by the application or by a middleware layer. PostgreSQL has table partitioning (splitting one logical table into physical child tables on the same server), which is related but different -- it does not give you multiple independent write paths.

\subsection{Range-Based Sharding}

In range-based sharding, rows are distributed by ranges of the sharding key. For example, users with IDs 1--10,000,000 go to shard 1, users 10,000,001--20,000,000 go to shard 2, and so on. Range sharding makes range queries efficient: ``find all users with IDs between 5,000,000 and 6,000,000'' hits exactly one shard.

The problem with range sharding is uneven load. If your newest users are the most active (a common pattern in social applications), shard $N$ (containing the highest user IDs) will be much hotter than shard 1. This is called a \emph{hotspot}.

\subsection{Hash-Based Sharding}

In hash-based sharding, a hash function is applied to the sharding key, and the result modulo the number of shards determines placement: \texttt{shard = hash(user\_id) \% num\_shards}. Hash sharding distributes load evenly across shards (assuming a good hash function and no key skew), eliminating hotspots.

The trade-off is that range queries become expensive. ``Find all orders for user 42'' is a single shard lookup. ``Find all orders created yesterday'' is a scatter-gather to all shards. Hash sharding optimizes for point lookups at the expense of range scans.

\begin{trapbox}{``I'll just add sharding from the start''}
Premature sharding is a classic over-engineering mistake. Sharding introduces enormous operational complexity: cross-shard joins become application-level code, foreign keys across shards cannot be enforced by the database, \texttt{ORDER BY} with \texttt{LIMIT} requires fetching from all shards and sorting in memory, and changing the shard count requires re-hashing and migrating massive amounts of data (resharding). Instagram ran its entire photo sharing service on a small number of PostgreSQL primaries for years before sharding, and spent enormous engineering effort sharding correctly when they actually needed it. Start with a well-indexed single primary with read replicas. Shard only when you have exhausted vertical scaling and read scaling, and only on the specific tables that are bottlenecks.
\end{trapbox}

\subsection{The Resharding Problem}

If you start with 8 shards and later need 16, you must move half the data from each existing shard to a new shard. During this migration, the system must continue serving traffic, which means routing logic is inconsistent, the migration might run for hours or days, and any error could leave data in the wrong shard. This is called resharding, and it is one of the hardest operational problems in distributed databases.

\begin{industrybox}{PlanetScale (Vitess) -- Online Schema Changes and Resharding}
PlanetScale is built on Vitess, the MySQL sharding middleware open-sourced by YouTube (Google) in 2012. Vitess solves resharding by treating it as a continuous process. It maintains a shard mapping table and can gradually migrate key ranges from one shard to another while serving traffic, using a shadow write approach (writes go to both old and new shard during migration, reads switch over once the new shard is consistent). PlanetScale's 2023 documentation describes their branching model as allowing schema changes to be tested on a branch without locking the production database -- they solve the ``ALTER TABLE locks the table'' problem by using Vitess's online schema change tool (gh-ost or pt-online-schema-change under the hood), which copies the table to a shadow copy, applies the schema change, and atomically swaps the tables.
\end{industrybox}

\newpage

% =============================================================================
\section{Connection Pooling: The Hidden Bottleneck}
% =============================================================================

PostgreSQL creates a new OS process for every database connection. A connection has a memory overhead of roughly 5--10 MB per connection (stack, shared memory segment, lock state). At 1,000 simultaneous connections, that is up to 10 GB of RAM consumed just by connection overhead, before any actual query work is done. More importantly, PostgreSQL's process-per-connection model means that each connection is a full OS process; context switching between 1,000 processes under load can consume significant CPU time.

The solution is connection pooling. A connection pool sits between the application and the database. The application has many lightweight threads or goroutines, each of which borrows a connection from the pool, executes a query, and returns the connection. The pool maintains a much smaller set of actual PostgreSQL connections.

\begin{termbox}{PgBouncer -- PostgreSQL Connection Pooler}
PgBouncer is the most widely used connection pooler for PostgreSQL. It runs as a lightweight daemon (typically on the same host as the application servers) and presents a PostgreSQL-compatible interface to application code. Behind the scenes, it multiplexes application connections onto a pool of actual server connections. PgBouncer's transaction-mode pooling allows a single server connection to be shared across multiple application connections, as long as each application connection is in ``between transactions'' state. Recommended pool size for PostgreSQL: the PostgreSQL documentation suggests that the number of active server connections should not exceed \texttt{2 * num\_CPU\_cores + effective\_spindle\_count}. For a 16-core database server, a pool of 30--40 server connections is typically sufficient to saturate the CPU without excess context switching.
\end{termbox}

\begin{industrybox}{Notion -- Connection Pooling at Scale (2023)}
Notion's engineering blog detailed their 2023 database scaling work. At their traffic peak, they were running hundreds of application server instances, each with its own connection pool. Without a centralized pooler, the total number of PostgreSQL connections grew proportionally to application server count, eventually hitting PostgreSQL's \texttt{max\_connections} limit. They deployed PgBouncer as a centralized pooler (rather than per-application-server), which reduced their effective server connection count from thousands to a few hundred, eliminating a significant category of ``too many connections'' errors during traffic spikes.
\end{industrybox}

\newpage

% =============================================================================
\section{Decision Framework: Relational vs. NoSQL}
% =============================================================================

The interview question ``what database would you use?'' is really several questions in disguise: What is the shape of the data? What are the access patterns? What consistency level does the application require? What is the write volume? Relational databases are the correct default for the vast majority of applications. NoSQL databases are the correct choice for specific, well-understood deviations from the relational model.

\begin{center}
\renewcommand{\arraystretch}{1.4}
\begin{tabularx}{\textwidth}{>{\bfseries}p{3.2cm} X X}
\toprule
\textbf{Dimension} & \textbf{Relational (PostgreSQL / MySQL)} & \textbf{NoSQL (various)} \\
\midrule
Data shape & Structured, tabular, normalized & Flexible: documents, key-value, wide-column, graph \\
\addlinespace
Relationships & First-class: foreign keys, joins & Application-level; no enforced referential integrity \\
\addlinespace
Consistency & ACID by default & Varies: eventual (Cassandra, DynamoDB) to strong (MongoDB 4+) \\
\addlinespace
Query flexibility & Arbitrary SQL, ad-hoc queries & Limited; queries optimized for specific access patterns \\
\addlinespace
Write scaling & Vertical + read replicas + sharding (complex) & Horizontal by design (Cassandra, DynamoDB) \\
\addlinespace
Schema changes & Requires migrations; \texttt{ALTER TABLE} can lock & Schema-free; flexible but no enforcement \\
\addlinespace
When to use & Transactional workloads, financial data, user data, anything with complex relationships & Massive write throughput, time-series, flexible schemas, graph traversals, session stores \\
\addlinespace
Common examples & User accounts, orders, inventory, payments & User sessions (Redis), activity feeds (Cassandra), product catalogs (MongoDB), analytics (ClickHouse) \\
\bottomrule
\end{tabularx}
\end{center}

The most important principle to internalize is that the relational vs. NoSQL decision is an \emph{access pattern decision}, not a performance or scale decision. Many teams have migrated from NoSQL back to PostgreSQL when they realized their data had more structure and relational integrity needs than they anticipated -- MongoDB's 2023 case studies include multiple enterprises that standardized on PostgreSQL after finding that ad-hoc schema flexibility created consistency bugs at scale.

\begin{trapbox}{``We need NoSQL for scale''}
This is a nearly universal misconception. PostgreSQL with proper indexing, connection pooling, and read replicas can handle tens of thousands of reads per second and thousands of writes per second on commodity hardware. Shopify's core commerce database ran on MySQL at millions of requests per day for years before sharding was necessary. The moment you encounter a scale argument for NoSQL, ask: have you exhausted read replicas? Have you profiled your slow queries? Have you checked your index coverage? The answer is almost always ``no.'' Relational databases can be scaled much further than most teams need before their limitations become relevant.
\end{trapbox}

\newpage

% =============================================================================
\section{Critical Numbers Reference}
% =============================================================================

\begin{center}
\renewcommand{\arraystretch}{1.5}
\begin{tabularx}{\textwidth}{>{\bfseries}p{5cm} X p{3.5cm}}
\toprule
\textbf{Metric} & \textbf{Value / Range} & \textbf{Source / Year} \\
\midrule
PostgreSQL connection overhead & 5--10 MB RAM per connection & PostgreSQL docs, 2024 \\
\addlinespace
Recommended active connections & $2 \times$ CPU cores + disk count & PostgreSQL wiki \\
\addlinespace
B-tree height for 10M rows & 4--5 levels & PostgreSQL internals docs \\
\addlinespace
Sequential scan speed (SSD) & $\sim$500 MB/s & Commodity NVMe SSD benchmark \\
\addlinespace
B-tree index lookup & 5--6 page reads ($O(\log n)$) & PostgreSQL internals \\
\addlinespace
Async replication lag (healthy) & $<$100 ms & PostgreSQL streaming replication \\
\addlinespace
Aurora storage write latency & $<$1 ms (quorum) & AWS Aurora blog, 2023 \\
\addlinespace
Aurora replication copies & 6 copies across 3 AZs & AWS Aurora architecture \\
\addlinespace
Aurora quorum write & 4 of 6 storage nodes & AWS Aurora blog, 2023 \\
\addlinespace
Shopify peak write throughput & $\sim$40,000 requests/sec (Black Friday 2023) & Shopify Engineering Blog, 2023 \\
\addlinespace
PgBouncer overhead & $\sim$2 KB RAM per client connection & PgBouncer documentation \\
\addlinespace
Typical \texttt{fsync()} latency (NVMe) & 0.1--0.5 ms & Linux kernel benchmarks \\
\addlinespace
PlanetScale (Vitess) max table size tested & $>$1 billion rows & PlanetScale docs, 2024 \\
\bottomrule
\end{tabularx}
\end{center}

\newpage

% =============================================================================
\section{System Design Application}
% =============================================================================

\subsection{URL Shortener (e.g., bit.ly)}

A URL shortener stores a mapping from a short code (7 characters) to the original URL. The read-to-write ratio is extreme: every URL is created once but clicked thousands or millions of times. The relational model fits perfectly: a single table \texttt{urls(code TEXT PRIMARY KEY, original\_url TEXT NOT NULL, created\_at TIMESTAMPTZ, user\_id BIGINT REFERENCES users(id))} with a B-tree primary key index on \texttt{code} gives $O(\log n)$ lookups. A covering index on \texttt{(code, original\_url)} eliminates the heap fetch for the redirect lookup (the most common operation).

The write path can go to a single PostgreSQL primary. The read path is the interesting scaling problem: at 100,000 redirects per second, even a fast index lookup cannot keep up with a single database. The solution is a caching layer (Redis) in front of the database -- the redirect lookup hits Redis first, and only falls back to PostgreSQL on a cache miss. The relational database becomes a source of truth and a cache-miss handler, not the primary serving path. The relational model was the right choice because the data is structured, the relationships (user owns URLs) are important, and SQL makes analytics queries (``how many clicks per URL per day?'') trivially easy.

\subsection{E-Commerce Order System (e.g., Shopify Merchant)}

An e-commerce order involves multiple related entities: orders, line items, customers, products, inventory. These relationships are inherently relational. A single order creation requires atomically inserting into \texttt{orders}, \texttt{order\_line\_items}, and decrementing \texttt{inventory}. If any step fails, none should persist. This is a textbook ACID transaction. Using NoSQL here would require the application to implement its own two-phase commit logic -- re-inventing what PostgreSQL already does correctly.

For read scaling, a read replica serves the product catalog, order history, and reporting queries. When a merchant's order table grows past 100 million rows, a \texttt{created\_at} index with range partitioning (PostgreSQL table partitioning by month) keeps index sizes manageable and allows old partitions to be archived to cheaper storage. Sharding by \texttt{merchant\_id} is the natural choice if write volume grows beyond a single primary, because almost all queries for a given merchant touch only that merchant's data.

\subsection{User Authentication Service}

An authentication service stores user credentials, sessions, and permissions. The data is highly relational (users have roles, roles have permissions), and the integrity requirements are extreme (two accounts must never share an email address, a session must always reference a valid user). A PostgreSQL unique constraint on \texttt{email} and a foreign key from \texttt{sessions.user\_id} to \texttt{users.id} enforce these invariants at the database layer -- they cannot be accidentally bypassed by a bug in the application code.

The access pattern for authentication is almost entirely point lookups: ``find user by email'' and ``find session by token.'' Both are served in $O(\log n)$ with B-tree indexes. At even high traffic (millions of logins per day), a single well-indexed PostgreSQL primary with PgBouncer can handle the load. The temptation to use Redis or a NoSQL store for ``simplicity'' typically backfires when the team realizes they need ACID semantics to safely handle concurrent password reset requests or session invalidation.

\newpage

% =============================================================================
\section{Curiosity Corner}
\addcontentsline{toc}{section}{Curiosity Corner}
% =============================================================================

\begin{curiobox}{What is MVCC doing with disk space, and why does VACUUM matter?}
Every \texttt{UPDATE} in PostgreSQL does not modify the existing row in place. Instead, it inserts a new version of the row (the updated tuple) and marks the old version as ``dead'' with a transaction ID stamp. The dead tuple stays on disk, invisible to new transactions but occupying space. Over time, in a write-heavy database, dead tuples accumulate and the table grows even if the logical data set is stable. \texttt{VACUUM} is the process that reclaims dead tuple space, making it available for new rows. \texttt{AUTOVACUUM} runs this automatically in the background. On a table receiving thousands of updates per second, a poorly tuned autovacuum can fall behind, causing ``table bloat'' -- the table grows to several times its logical size. Notion documented in 2023 that table bloat was causing their sequential scans to read far more data than the table logically contained, degrading performance. The fix was aggressive autovacuum tuning: lower \texttt{autovacuum\_vacuum\_scale\_factor} so vacuum triggers earlier, and increase \texttt{autovacuum\_vacuum\_cost\_delay} budget so it runs faster. Monitoring table bloat and dead tuple ratio is part of PostgreSQL operations at scale.
\end{curiobox}

\begin{curiobox}{What is the difference between PostgreSQL table partitioning and sharding?}
PostgreSQL table partitioning (introduced as declarative partitioning in version 10, 2017) splits one logical table into multiple physical child tables on the \emph{same database server}. Queries against the parent table are automatically routed to the relevant child tables (partition pruning). This improves performance by reducing the size of each individual index and allowing old partitions to be dropped or archived cheaply. It does not give you multiple independent write paths -- all partitions still share the same single primary server. Sharding, by contrast, distributes data across multiple \emph{independent} database servers, each with its own primary and write capacity. Partitioning is a local optimization; sharding is a distributed systems decision. Use partitioning first (it is much simpler), and shard only when partitioning is insufficient.
\end{curiobox}

\begin{curiobox}{What is a WAL and why does it get streamed to replicas?}
The Write-Ahead Log (WAL) is a sequential log of every change made to the database, written to disk before the change is applied to the data pages. It exists for two reasons: crash recovery (replay the log on startup after a crash to restore consistency) and replication (stream the log to replicas so they can apply the same changes). The WAL is a sequential binary log, which means writing to it is much faster than random writes to data pages -- a huge advantage on spinning disks, and still meaningful on SSDs because sequential I/O saturates higher bandwidth. Replicas receive WAL records over a TCP connection (streaming replication) and apply them continuously. The WAL is also used by logical replication (PostgreSQL 10+), which replays changes at the SQL statement level rather than the binary page level, allowing replication to heterogeneous targets (e.g., replicating from PostgreSQL to Kafka via pgoutput or wal2json plugins).
\end{curiobox}

\begin{curiobox}{How does ``SELECT FOR UPDATE'' prevent double-spending in a payments system?}
Consider a naive implementation of a payment deduction: read the account balance, check if it is sufficient, then write the decremented balance. If two requests run concurrently (two charges of \$50 to an account with \$60), both might read the balance as \$60, both pass the sufficiency check, and both write \$10 -- leaving the account at \$10 when it should have been declined for the second charge. This is a classic race condition. \texttt{SELECT ... FOR UPDATE} acquires an exclusive row-level lock on the account row at read time. The second transaction's \texttt{SELECT FOR UPDATE} blocks until the first transaction commits. When it resumes, it re-reads the now-\$10 balance and correctly declines the charge. The sequence of operations -- lock, read, check, write, commit -- is now serialized for this row, eliminating the race. This is why pessimistic locking is appropriate for financial transactions, even though it is slower than optimistic concurrency: the cost of a false positive (overdraft) far exceeds the cost of serialized access.
\end{curiobox}

\begin{curiobox}{Why is the N+1 query problem a database-level concern, not just an ORM problem?}
The N+1 query problem occurs when an application executes one query to fetch $N$ parent records, then executes $N$ separate queries to fetch the related child records for each parent. For 100 blog posts with their authors, this is 1 + 100 = 101 queries. ORMs like Rails ActiveRecord and Django ORM make this problem very easy to create accidentally, but the root cause is at the database connection level: 101 round trips to PostgreSQL, each incurring network latency and connection overhead. The fix -- a \texttt{JOIN} query or an \texttt{IN} clause to fetch all children in one query -- reduces 101 queries to 1 or 2. At scale, N+1 patterns have killed production databases: 10 requests per second times 101 queries each is 1,010 queries per second instead of 10--20. GitHub's engineering blog has documented several incidents where N+1 patterns introduced during code review caused database CPU spikes that affected production traffic. Recognizing and fixing N+1 patterns is a senior engineering competency.
\end{curiobox}

\begin{curiobox}{What is an XID wraparound and why can it bring down a PostgreSQL database?}
PostgreSQL uses 32-bit transaction IDs (XIDs) to track which transactions created which row versions (for MVCC). With a 32-bit counter, XIDs cycle after approximately 2 billion transactions (about $2^{31}$, since one bit is reserved). When XIDs wrap around, PostgreSQL could confuse ``future'' transaction IDs with ``past'' ones, making rows appear invisible. To prevent this catastrophe, PostgreSQL requires VACUUM to freeze old row versions -- replacing their XID with a special ``frozen'' flag that is always visible -- before the wraparound point. If a database is not vacuumed frequently enough and XID exhaustion approaches (within 3 million transactions by default), PostgreSQL enters a safety shutdown mode, refusing new connections and printing a message saying ``database is not accepting commands to avoid wraparound data loss.'' This is not hypothetical: in 2019, a Mailchimp PostgreSQL database hit XID wraparound due to a stalled autovacuum, taking the service offline until emergency VACUUM could be run. Monitoring XID age is a non-negotiable operational requirement for any PostgreSQL deployment.
\end{curiobox}

\newpage

% =============================================================================
\section{Self-Quiz}
\addcontentsline{toc}{section}{Self-Quiz}
% =============================================================================

\begin{quizbox}
\textbf{Answer these as an interviewer would probe you -- not with definitions, but with reasoning and tradeoffs.}

\vspace{6pt}
\textbf{Q1.} You are designing a payment processing service. The product manager says ``we should use MongoDB because it's more flexible and faster.'' How do you respond, and what do you actually choose?

\vspace{10pt}
\textbf{Q2.} You have a \texttt{messages} table with 500 million rows. A query \texttt{SELECT * FROM messages WHERE sender\_id = 42 ORDER BY created\_at DESC LIMIT 20} is taking 800 ms. Walk me through your diagnostic process and fix.

\vspace{10pt}
\textbf{Q3.} Your PostgreSQL primary is handling 50,000 reads per second and is at 90\% CPU. Your operations team proposes adding a second primary. What's wrong with that idea, and what do you do instead?

\vspace{10pt}
\textbf{Q4.} In a read committed isolation level transaction, you run the same \texttt{SELECT} query twice. Can you get different results? Why or why not, and when would this matter in a real system?

\vspace{10pt}
\textbf{Q5.} You are designing the database layer for a multi-tenant SaaS application where each customer has their own isolated dataset. Describe at least two sharding strategies and their tradeoffs for this specific scenario.
\end{quizbox}

\newpage

% =============================================================================
\section*{Model Answers}
\addcontentsline{toc}{section}{Model Answers}
% =============================================================================

\begin{answerbox}
\textbf{A1. Payment processing: ACID is non-negotiable.}

Push back directly: flexibility is not a database requirement for payments -- correctness is. A payment deduction requires atomically updating the account balance and recording the transaction ledger entry. If those two operations are not atomic, you will have ledger entries without balance changes (money appears from nowhere) or balance changes without ledger entries (money disappears without record). MongoDB has supported multi-document ACID transactions since version 4.0 (2018), so the ``MongoDB can't do ACID'' argument is outdated -- but the document model is a poor fit for relational financial data. Choose PostgreSQL. Use \texttt{SELECT FOR UPDATE} for balance checks, wrap balance and ledger writes in a single transaction, and use read replicas for reporting queries. The ``faster'' claim needs measurement, not assumption.

\vspace{8pt}
\textbf{A2. Diagnosing the slow query on messages.}

Start with \texttt{EXPLAIN (ANALYZE, BUFFERS)} to see the query plan. Likely finding: the query is doing a sequential scan (no usable index) or an index scan on \texttt{sender\_id} followed by a sort (no covering index including \texttt{created\_at}). Fix: create a composite index on \texttt{(sender\_id, created\_at DESC)}. With this index, PostgreSQL can scan the index in reverse order for \texttt{sender\_id = 42}, reading exactly 20 rows without a separate sort step. If the index already exists and is not being used, check whether the query planner is choosing a sequential scan because the index statistics are stale (\texttt{ANALYZE messages}) or because \texttt{sender\_id} has very low selectivity (one user has half the messages). In the low-selectivity case, consider whether a partial index or application-level archival is appropriate.

\vspace{8pt}
\textbf{A3. No such thing as a second primary -- use read replicas.}

A ``second primary'' in traditional replication does not exist without conflict resolution logic. Both primaries would accept writes, and without coordination, the same row could be updated differently on each primary (split-brain). The correct answer is to add read replicas: if 90\% of your load is reads (a common ratio for consumer applications), three read replicas reduce primary read load by 75\%. Deploy PgBouncer to manage connections. Profile the remaining primary CPU: is it from query execution or from write WAL overhead? If writes are the CPU bottleneck, the next step is query optimization and potentially sharding the specific hot tables.

\vspace{8pt}
\textbf{A4. Non-repeatable reads in Read Committed.}

Yes, you can get different results. In Read Committed isolation, each statement within a transaction sees a fresh snapshot of committed data -- not the snapshot at transaction start. If another transaction commits a change to a row between your two \texttt{SELECT} statements, your second \texttt{SELECT} will see the new value. This matters in banking: if your transaction reads an account balance, another transaction commits a debit, and your transaction reads again, you now have an inconsistent view of the balance. The fix is to use Repeatable Read isolation for transactions that need a consistent view of the same data across multiple statements, or to use \texttt{SELECT FOR UPDATE} to lock the rows you are reading.

\vspace{8pt}
\textbf{A5. Sharding strategies for multi-tenant SaaS.}

The natural sharding key is \texttt{tenant\_id} (also called \texttt{customer\_id} or \texttt{org\_id}). Two strategies:

Strategy 1 -- \textbf{Hash sharding on tenant\_id}: \texttt{shard = hash(tenant\_id) \% N}. Even distribution of tenants across shards. All queries for a single tenant (the dominant pattern) hit exactly one shard. Cross-tenant analytics requires scatter-gather. Does not isolate large ``whale'' tenants that generate disproportionate load. Good for: SaaS with many similarly-sized tenants.

Strategy 2 -- \textbf{Dedicated shard per tenant tier}: Small/free tenants share a few high-density shards. Enterprise tenants get dedicated shards (or even dedicated clusters). This is how Shopify partitions merchant data -- small merchants share infrastructure, large merchants (e.g., Kylie Cosmetics) get dedicated resources. More operationally complex but eliminates noisy-neighbor problems and allows SLA differentiation by tier.

For both strategies: route queries through a shard map (a small relational database mapping tenant to shard). The shard map itself is tiny and can be aggressively cached.
\end{answerbox}

\newpage

% =============================================================================
\section{What's Next: L1-06 -- Databases (NoSQL)}
\addcontentsline{toc}{section}{What's Next}
% =============================================================================

\begin{insightbox}
\textbf{You now understand what relational databases guarantee and how they deliver those guarantees.} The ACID contract, the B-tree index structure, the WAL, MVCC, primary-replica replication, and the sharding decision framework -- these are the foundations that every serious database conversation builds on. You can now reason precisely about when PostgreSQL is the right choice and what its limits are, rather than guessing.

The next lesson, \textbf{L1-06 -- Databases: NoSQL}, inverts the frame. Rather than asking ``how does PostgreSQL guarantee consistency?'', L1-06 asks: ``what do you \emph{give up} when you need write throughput that a single primary cannot provide?'' You will study the CAP theorem and what it means in practice (not as a theoretical trilemma but as a concrete engineering choice), Cassandra's wide-column model and how it achieves linear write scaling at the cost of no joins and tunable consistency, DynamoDB's single-table design pattern, and the Redis data model for ephemeral in-memory workloads.

The sequence that leads to your first mock interview: L1-06 NoSQL $\rightarrow$ L1-07 Caching $\rightarrow$ L1-08 Message Queues, which together unlock Gate 1 and the Notification System mock (L3-01). Every concept you just learned -- replication lag, consistency guarantees, connection overhead -- will appear again in those lessons as contrasting points. The relational database is the baseline; everything else is a deviation from it, justified by specific constraints.
\end{insightbox}

\vspace{1.5cm}
\begin{center}
  {\color{highlight}\rule{0.5\linewidth}{0.6pt}}

  \vspace{0.4cm}
  {\small\color{primary!60}
    L1-05 --- Databases: Relational \quad|\quad System Design Interview Prep\\
    Edition 2025 \quad|\quad Compiled with \texttt{pdflatex}
  }
\end{center}

\end{document}
